\chapter{Future Work: Specialized Indexes on Graph Data}
\label{chap:specialized_indexing}

Indexes. Most graph database systems use indexes. Now, systems based on non-graph backends, for example RDBMS or document stores, usually rely on existing indexing infrastructure present in such systems. Native graph databases employ index structures for the neighborhoods of each vertex, often in the form of direct pointers [175]. Neighborhood indexes are used mostly to speed up the access of adjacency lists to accelerate traversal queries. JanusGraph calls these indexes vertex centric. They are constructed specifically for vertices, so that incident edges can be filtered efficiently to match the traversal conditions [21]. While JanusGraph allows multiple vertex-centric indexes per vertex, each optimized for different conditions, which are then chosen by the query optimizer, simpler solutions exist as well. LiveGraph uses a two-level hierarchy, where the first level distinguishes edges by their label, before pointing to the actual physical storage [212]. Graphflow indexes the neighbors of a vertex into forward and backward adjacency lists, where each list is first partitioned by the edge label and, second, by the label of the neighbor vertex [120]. Another example is Sparksee, which uses various different index structures to find the adjacent vertices and properties of a vertex [139]. An interesting research opportunity is to design indexes for richer “higher-order” structural information beyond plain neighborhoods. Specifically, a recent wave of pattern matching schemes [35, 146] indicates the importance of higher-order graph structure, such as triangles to which each vertex belongs. Indexing such information would significantly speed up queries related to clique mining, dense subgraph discovery, clustering, and many others. While some work has been done in this respect [182], many designs could be proposed. Data indexes concern data beyond the neighborhood information, and they can be used to accelerate query plan generation and query execution. It is possible, for example, to index all vertices that have a specific property (value). They are usually employed to speed up Business Intelligence workloads (details on workloads are in Section 4). Many triple stores, for example AllegroGraph [82], provide all six permutations of subject (S), predicate (P), and object (O) as well as additional aggregated indexes. However, to reduce associated costs, other approaches exist as well: TripleBit uses just two permutations (PSO, POS) with two aggregated indexes (SP, SO) and two auxiliary index structures [210]. gStore implements pattern matching queries with the help of two index structures: a VS*-tree, which is a specialized B+-tree, and a trie-based T-index [213]. Some database systems like Amazon Neptune [5] or AnzoGraph [48] only provide implicit indexes, while still being confident to answer all kinds of queries efficiently. However, most graph database systems allow the user to explicitly define data indexes. Some of them, like Azure Cosmos DB [148], support composite indexes (a combination of different labels/properties) for more specific use cases. In addition to internal indexes, some systems employ external indexing tools. For example, Titan and JanusGraph [21] use internal indexing for label- and value-based lookups but rely on external indexing backends (e.g., Elasticsearch [72]orApacheSolr[15]) for non-trivial lookups involving multiple properties, ranges, or full-text search. Structural indexes are used for various internal data. Here LiveGraph uses a vertex index to map its vertex IDs to a physical storage location [212]. ArangoDB uses a hybrid index, a hashtable, to find the documents of incident edges and adjacent vertices of a vertex [16]. We categorize systems (for which we were able to find this information) according to this criteria in Table 5. We find no clear connection between the index type and the backend of a graph database, but most systems use tree-based indexes. A research opportunity, useful especially for practitioners, would be to conduct a detailed and broad performance comparison of different index implementations for different workloads.